\subsection{整代数的维数}

引理 2. --- 设 \(\rho: A \to B\) 是一个环同态，使得对 A 的每个素理想 \(\mathfrak{p}\)，\(\kappa(\mathfrak{p})\)-代数 \(B \otimes_{A} \kappa(\mathfrak{p})\) 是整的（V, § 1, no. 1, 定义 2）。设 \(\mathfrak{q}\) 和 \(\mathfrak{q}'\) 是 B 的两个素理想，满足 \(\mathfrak{q} \subset \mathfrak{q}'\) 且 \(\mathfrak{q} \neq \mathfrak{q}'\)。则 \(\rho^{-1}(\mathfrak{q}) \neq \rho^{-1}(\mathfrak{q}')\)。

事实上，如果 \(\mathfrak{q}\) 和 \(\mathfrak{q}'\) 位于 A 的同一个素理想 \(\mathfrak{p}\) 之上，则根据 no. 1 的引理 1，有 \(\dim(B \otimes_{A} \kappa(\mathfrak{p})) \geqslant 1\)，这与 \(\dim(B \otimes_{A} \kappa(\mathfrak{p})) \leqslant 0\) 这一事实矛盾（§ 1, no. 3, 例 6）。

定理 1. --- 设 \(\rho: A \to B\) 是一个环同态，使 \(B\) 成为整 \(A\)-代数。

\begin{enumerate}
\def\labelenumi{\alph{enumi})}
\item
  设 M 是一个有限生成 A-模。则有 \(\dim_{\mathbf{B}}(\mathbf{M} \otimes_{\mathbf{A}} \mathbf{B}) \leqslant \dim_{\mathbf{A}}(\mathbf{M})\)。特别地，有 \(\dim(\mathbf{B}) \leqslant \dim(\mathbf{A})\)。如果映射 \({}^{a}\rho: \operatorname{Spec}(\mathbf{B}) \to \operatorname{Spec}(\mathbf{A})\) 是满射，例如当 \(\rho\) 为单射时（V, § 2, no. 1, 定理 1），则有 \(\dim_{\mathbf{B}}(\mathbf{M} \otimes_{\mathbf{A}} \mathbf{B}) = \dim_{\mathbf{A}}(\mathbf{M})\)，特别地有 \(\dim(\mathbf{B}) = \dim(\mathbf{A})\)。
\item
  设  是 B 的一个理想，且 \(\mathfrak{a} = \rho^{-1}(\mathfrak{b})\) 是它在 A 中的逆像。我们有 \(\operatorname{ht}(\mathfrak{b}) \leqslant \operatorname{ht}(\mathfrak{a})\) 且 \(\dim(\mathbf{B}/\mathfrak{b}) = \dim(\mathbf{A}/\mathfrak{a})\)。如果 \({}^{a}\rho : \operatorname{Spec}(B) \to \operatorname{Spec}(A)\) 是满射，则对 A 的每个理想  ，有 \(\operatorname{ht}(\mathfrak{a}\mathbf{B}) \leqslant \operatorname{ht}(\mathfrak{a})\)。
\item
  假设 B 在 A 上有限，且 N 是一个有限生成 B-模。则有 \(\dim_{\mathbf{B}}(\mathbf{N}) = \dim_{\mathbf{A}}(\mathbf{N})\)。特别地，有 \(\dim(\mathbf{B}) = \dim_{\mathbf{A}}(\mathbf{B})\)。
\end{enumerate}

证明 a)。根据 V, § 1, no. 1 的命题 5，\(\kappa(\mathfrak{p})\)-代数对于 A 的每个素理想 \(\otimes_{\mathbf{A}} \kappa(p)\)\(\mathfrak{p}\) 都是整的。设 \(\mathfrak{q}_0 \subset \ldots \subset \mathfrak{q}_m\) 是 B 的素理想链；根据引理 2，理想 \(\mathfrak{p}_i = \rho^{-1}(\mathfrak{q}_i)\) 两两不同，因而 \(\mathfrak{p}_0 \subset \ldots \subset \mathfrak{p}_m\) 是 A 的素理想链，从而 \(m \leqslant \dim(\mathbf{A})\)。因此 \(\dim(\mathbf{B}) \leqslant \dim(\mathbf{A})\)。

现在假设映射 \({}^{a}\rho\) 是满射。设 \(\mathfrak{p}_0 \subset \ldots \subset \mathfrak{p}_n\) 是 A 的素理想链；于是存在一个素理想 \(\mathfrak{q}_0\)，位于 \(\mathfrak{p}_0\) 之上。根据第一存在定理的推论 2（V, § 2, no. 1, th. 1），通过对 \(n\) 作归纳，可以构造 B 的素理想链 \(\mathfrak{q}_0 \subset \ldots \subset \mathfrak{q}_n\)，使得 \(\mathfrak{q}_i\) 位于 \(\mathfrak{p}_i\) 之上，对 \(0 \leqslant i \leqslant n\) 成立。因此 \(n \leqslant \dim(\mathbf{B})\)，从而 \(\dim(\mathbf{A}) \leqslant \dim(\mathbf{B})\)。

这证明了在断言 \(a\) 的 \(\mathbf{M} = \mathbf{A}\) 情形下成立。在一般情形中，令 \(\alpha\) 为 \(\mathbf{M}\) 的湮灭子，于是 \(\mathbf{M}\) 的支集可识别为 ，并有 \(\operatorname{Spec}(\mathbf{A}/\alpha)\)\(\dim_{\mathbf{A}}(\mathbf{M}) = \dim(\mathbf{A}/\alpha)\)。根据 II, § 4, no. 4 的命题 19，\(\mathbf{M} \otimes_{\mathbf{A}} \mathbf{B}\) 的支集是通过 \(^{\alpha}\rho\) 得到的 \(\mathbf{M}\) 的支集的逆像，因而可识别为 ，并有 \(\operatorname{Spec}(\mathbf{B}/\rho(\alpha)\mathbf{B})\)。只需注意，同态 \(\dim_{\mathbf{B}}(\mathbf{M} \otimes_{\mathbf{A}} \mathbf{B}) = \dim(\mathbf{B}/\rho(\alpha)\mathbf{B})\)\(\rho': A/\alpha \to B/\rho(\alpha)\) 由 \(\rho\) 诱导，它使  成为整的 \(B/\rho(\alpha)\)-代数，并且 \((A/\alpha)\) 在 \(^{\alpha}\rho'\) 满射时，\(^{\alpha}\rho\) 也是满射的。

证明 \(b\)。根据 § 1, no. 3 的命题 7，只需证明，对于 A 的每个包含  的素理想  ，都有 ；设  为这样的理想。根据 V, § 2, no. 1, 定理 1 的推论 2，存在 B 的一个位于  之上且包含  的素理想 \_，并且根据 § 1, no. 3 的命题 7 有 \_。

现在 \(\mathbf{B}_{\mathfrak{q}}\) 可识别为整的 \(\mathbf{A}_{\mathfrak{p}}\)-代数 \(\mathbf{B} \otimes_{\mathbf{A}} \mathbf{A}_{\mathfrak{p}}\) 的一个分式环，从而

\[
\dim (\mathrm{B} _ {\mathfrak {q}}) \leqslant \dim (\mathrm{B} \otimes_ {\mathrm{A}} \mathrm{A} _ {\mathfrak {p}}) \leqslant \dim (\mathrm{A} _ {\mathfrak {p}})
\]

根据 § 1, no. 3 的命题 6 及上述断言 a)，我们由此证明了不等式 。此外，由  诱导的  到  的同态是单射，并使  成为整的 -代数；因此由 a) 有  。假设  是满射，设  是 A 的一个理想， 是 A 的一个包含  的素理想。根据假设，存在 B 的一个位于  之上的素理想  。由前述结果有  ，从而  。取下确界，得到  。

最后，  由应用于 N 的湮灭子  的  得出。

定理 2. --- 设 A 是一个整闭环，B 是包含 A 且在 A 上整的环。假设 B 是无挠 A-模。对 A 的每个理想  ，有 。设  是 B 的一个理想且 ；则有 。

设 \(\rho\) 是从 A 到 B 的典范映射。设 \(\mathfrak{a}\) 是 A 的一个理想。如果 \(\mathfrak{a} = \mathbf{A}\)，第一个等式显然成立。设 \(\mathfrak{a} \neq \mathbf{A}\)。由于 \(\rho\) 是单射，\({}^{a}\rho\) 是满射（V, § 2, no. 1, 定理 1）。因此 \(\mathfrak{aB} \neq \mathbf{B}\)。现在设 \(\mathfrak{q}\) 是 B 的一个包含 \(\mathfrak{aB}\) 的素理想。置 \(\mathfrak{p} = \mathfrak{q} \cap \mathbf{A}\)。有 \(\mathfrak{a} \subset \mathfrak{p}\)，从而有 。设 \(\mathfrak{p}_0 \subset \ldots \subset \mathfrak{p}_n\) 是 A 的一条满足 \(\mathfrak{p}_n = \mathfrak{p}\) 的素理想链。根据第二存在定理（V, § 2, no. 4, th. 3），通过归纳可以构造 B 的一条素理想链 \(\mathfrak{q}_0 \subset \ldots \subset \mathfrak{q}_n\)，使得 \(\mathfrak{q}_n = \mathfrak{q}\)，且 \(\mathfrak{q}_i\) 位于 \(\mathfrak{p}_i\) 之上，对 \(0 \leqslant i \leqslant n\) 成立。有  ，从而有  。取下确界，得到  （§ 1, no. 3, 命题 7）。不等式  由定理 1 得出，从而得到第一个等式。设 \(n \leqslant\) 是 B 的一个理想。置 \(\mathfrak{a} = \rho^{-1}(\mathfrak{b})\)。有  ，从而有  。不等式  由定理 1 得出，从而得到该定理。

注。--- 设 A 是一个整环，B 是包含 A 且在 A 上整的环。设  是 A 的一个素理想，使得  的整闭包是一个局部环。可以证明，当 B 是整环时，对于 B 中位于  之上的每个素理想  ，有 \(A_{p}\)（p. 85, 练习 9）。

